% source: https://tex.stackexchange.com/a/765836 (question 643403, score 4, block 1)
\documentclass{standalone}
\usepackage{fourier-otf}
\usepackage[svgnames]{xcolor}
\usepackage[3d]{luadraw}
\begin{document}
%https://tex.stackexchange.com/questions/643403/how-to-draw-a-shaded-boundary-using-tikz
\begin{luadraw}{name=shaded_boundary}
local ld = luadraw
local M, vecI, vecJ = ld.pt3d.M, ld.pt3d.vecI, ld.pt3d.vecJ
local g = ld.graph3d:new{ window3d={0,5,0,5,0,5}, adjust2d=true, viewdir={-20,70}, size={10,10} }
local a = -2
local k = function(u,v)
    local N = math.sqrt(1+v*(a+v))/v
    return M(1/v + N*math.exp(u), 1/v + N*math.exp(-u), v)
end
local S = ld.surface(k,-5,5,1e-4,5,{60,30})
S = ld.cutfacet(S, {M(5,0,0),-vecI})
S = ld.cutfacet(S, {M(0,5,0),-vecJ})
local Sv, Sh = g:Classifyfacet(S)
local Fx = ld.parametric3d( function(t) return M(5,(t+5+a)/(5*t-1),t) end, (10+a)/24,5)[1]
local Fy = ld.parametric3d( function(t) return M((t+5+a)/(5*t-1),5,t) end, (10+a)/24,5)[1]
local Fz = ld.parametric3d( function(t) return M(t,(t+5+a)/(5*t-1),5) end, (10+a)/24,5)[1]
table.insert(Fx, M(5,5,5)); table.insert(Fy, M(5,5,5)); table.insert(Fz, M(5,5,5))
g:Dboxaxes3d({grid=true, gridcolor="gray", drawbox=true,zlegendsep=0.25})
g:Fillopacity(0.8)
g:Dpolyline3d( Fy, true,"left color=green!20, right color=green")
g:Dpolyline3d( ld.border(Sh), "left color=cyan!20, right color=cyan")
g:Dpolyline3d( ld.border(Sv), "left color=blue!30, right color=blue, middle color=blue!25")
g:Dpolyline3d( Fx, true,"left color=yellow!20, right color=yellow")
g:Dpolyline3d( Fz, true,"left color=red!20, right color=red")
g:Show()
\end{luadraw}
\end{document}
